\section{Contexto y desafíos del curso}

Esta clase de UC Berkeley CS294-196, otoño de 2024, la impartió Yuandong Tian, de Meta AI FAIR, con el título ``Towards a Unified Framework of Neural and Symbolic Decision Making''; se centra en cómo unificar las redes neuronales, la búsqueda, la planeación y el ciclo cerrado de datos de agentes en un mismo marco de decisión.
Abrió con una tendencia progresiva: a medida que aumenta la dificultad de la tarea, la curva de desempeño del mejor modelo sigue bajando de manera monótona, sin mostrar elasticidad.
\textit{``This is actually in some sense a very worrisome kind of curve that even the best model has not solved this planning problem.''}
Al observar repetidamente varios benchmarks (planeación inteligente, el juego Soang, la búsqueda de rutas en Maze), la misma tendencia aparece una y otra vez: la puntuación del modelo queda acotada por las propiedades simbólicas implícitas en el grafo de cómputo, y simplemente apilar parámetros no permite abrir el panorama en poco tiempo.
Esto marca el tono de las tres vías de respuesta de esta clase: depender solo de modelos más grandes y más costosos no basta.

\subsection{La curva en la capacidad de planeación}

En los primeros 10 minutos, el ponente mostró una curva de evaluación para tareas de planeación complejas: aunque la precisión del modelo puede mejorar, la puntuación tras la convergencia sigue siendo baja, y la cantidad de tokens o parámetros de entrenamiento necesarios crece de manera exponencial.
Este fenómeno se repite en diálogo, planeación automática y navegación robótica; la raíz del problema está en que un modelo generativo probabilístico no puede garantizar propiedades simbólicas (restricciones, factibilidad).
\begin{knowledgebox}{Síntomas del fracaso en la planeación}
Los síntomas más comunes de los LLM en la planeación incluyen: 1) violaciones de restricciones (por ejemplo, ignorar límites o restricciones de consumo de energía); 2) acumulación de errores en horizontes largos; 3) ausencia de una gestión precisa de las bifurcaciones del espacio de estados. Cualquier problema que dependa de la precisión, como la programación logística o la aprobación de cumplimiento normativo, se ve amplificado por esto.
\end{knowledgebox}

\subsection{Tres estrategias}

El ponente propuso tres vías de corrección: \textbf{Scaling} (aumentar tokens/parámetros), \textbf{Hybrid} (combinar LLM y solver) y, de manera más sistemática, \textbf{Compound AI}. No son opciones mutuamente excluyentes, sino una arquitectura organizada por capas según el costo y el alcance del razonamiento, y se pueden apilar una sobre otra.
\begin{enumerate}
    \item El extremo de \textbf{Scaling} confía en la generalización probabilística mediante un modelado del lenguaje más grande, más datos de entrenamiento y más cómputo, pero es poco sensible al conocimiento estructural de las variables de restricción.
    \item El extremo de \textbf{Hybrid} usa el LLM para aportar candidatos/acciones más heurísticas, mientras el solver se encarga de hacer cumplir las restricciones, reduciendo la falta de control del modelo generativo puro.
    \item El extremo de \textbf{Compound AI} encadena varios módulos (traza de búsqueda, diálogo de múltiples turnos, solver con restricciones) para formar un ciclo de retroalimentación auditable, lo que facilita la división del trabajo y la depuración.
\end{enumerate}
\begin{figure}[H]
\centering
\includegraphics[width=\textwidth]{frame-solution-categories.jpg}
\caption{El ponente enumera tres categorías de estrategias de solución, ubicadas en distintos niveles de costo de ingeniería y alcance de razonamiento. \protect\footnotemark}
\end{figure}
\footnotetext{Intervalo de tiempo en el video: 00:09:10--00:09:18.}
\begin{importantbox}{Desglose de las tres estrategias}
Scaling apuesta fuertemente por los recursos, Hybrid deja que el LLM aporte heurísticas mientras el solver se encarga de la búsqueda, y Compound AI encadena varios módulos para formar un ciclo de retroalimentación integrado, lo cual permite tanto controlar la calidad como facilitar el diagnóstico.
\end{importantbox}
\begin{warningbox}{No confíe ciegamente en Scaling}
Depender exclusivamente de modelos más grandes implica ciclos de entrenamiento más largos, sobreajuste más difícil de depurar y un consumo de energía en inferencia excesivo; si no se observa una mejora estructural, los parámetros adicionales terminarán únicamente por amplificar errores que ``parecen razonables''.
\end{warningbox}

\subsection{Resumen de la sección}

El mecanismo generativo de los LLM presenta una incertidumbre inherente en los problemas de planeación, y se necesita recurrir a estructuras simbólicas para recuperar el control; las tres estrategias ofrecen distintos puntos de entrada, tanto de ingeniería como metodológicos.

